\chapter{Spinor Fields and the Dirac Equation}

\begin{remarkbox}
Spinor fields are the natural relativistic fields for spin-one-half degrees of
freedom. They are not ordinary scalars, vectors, or tensors. Their transformation
law is tied to the double cover of the Lorentz group, and their dynamics is
governed by a first-order relativistic equation: the Dirac equation.
\end{remarkbox}

\section{Why Spinors Are Needed}

Scalar fields describe spin-zero degrees of freedom. Vector fields describe
spin-one degrees of freedom, as in electromagnetism. However, many fundamental
matter fields are spin-one-half fields. Such fields cannot be represented by
ordinary spacetime scalars or vectors. They require spinors.

The need for spinors can be seen from symmetry. Relativistic fields must carry
representations of the Lorentz group. Scalars transform trivially, vectors
transform with Lorentz matrices \(\Lambda^\mu{}_{\nu}\), and tensors transform
with tensor products of these matrices. Spinors transform under a different kind
of representation: a representation of the double cover of the proper Lorentz
group.

The proper orthochronous Lorentz group \(SO^+(1,3)\) has a double cover,
\[
    SL(2,\mathbb{C})\longrightarrow SO^+(1,3).
\]
Spinors transform naturally under \(SL(2,\mathbb{C})\). This is why a rotation by
\(2\pi\) can act as multiplication by \(-1\) on a spinor, while a rotation by
\(4\pi\) returns it to itself.

\begin{definitionbox}
A \emph{spinor} is a field transforming under a spin representation of the
Lorentz group, more precisely under a representation of the double cover of the
proper Lorentz group.
\end{definitionbox}

In classical field theory, one may first treat spinors formally as fields with
components. The deeper quantum statement is that spin-one-half fields must be
quantized with anticommutation relations. This chapter focuses on the classical
field equations and their symmetries, while indicating where quantum field
theory enters.

\section{Clifford Algebra and Gamma Matrices}

The Dirac equation is built from gamma matrices. With the mostly-minus metric
convention
\[
    \eta_{\mu\nu}=\operatorname{diag}(+1,-1,-1,-1),
\]
the gamma matrices satisfy the Clifford algebra
\[
    \{\gamma^\mu,\gamma^\nu\}
    =\gamma^\mu\gamma^\nu+\gamma^\nu\gamma^\mu
    =2\eta^{\mu\nu}\mathbb{I}.
\]
Thus
\[
    (\gamma^0)^2=\mathbb{I},
    \qquad
    (\gamma^i)^2=-\mathbb{I}
\]
for spatial indices \(i=1,2,3\).

The gamma matrices are not themselves spacetime vectors. They are constant
matrices acting on spinor components, arranged so that contractions such as
\[
    \gamma^\mu\partial_\mu
\]
transform covariantly when the spinor transforms appropriately.

\begin{definitionbox}
The \emph{Clifford algebra} associated with Minkowski spacetime is generated by
matrices \(\gamma^\mu\) satisfying
\[
    \{\gamma^\mu,\gamma^\nu\}=2\eta^{\mu\nu}\mathbb{I}.
\]
\end{definitionbox}

A Dirac spinor in four spacetime dimensions has four complex components:
\[
    \psi(x)=
    \begin{pmatrix}
    \psi_1(x)\\
    \psi_2(x)\\
    \psi_3(x)\\
    \psi_4(x)
    \end{pmatrix}.
\]
Different explicit representations of the gamma matrices are possible. Physical
statements do not depend on the representation. A change of gamma-matrix basis
is a similarity transformation.

The key algebraic point is that the Clifford algebra linearizes the
Klein--Gordon operator. If \(\psi\) satisfies the Dirac equation
\[
    (i\gamma^\mu\partial_\mu-m)\psi=0,
\]
then applying \((i\gamma^\nu\partial_\nu+m)\) gives
\[
    (\Box+m^2)\psi=0.
\]
Thus every component of a free Dirac spinor satisfies the Klein--Gordon equation,
but the Dirac equation contains additional spinor structure.

\section{Lorentz Transformations of Spinors}

Under a Lorentz transformation
\[
    x^\mu\longrightarrow x'^\mu=\Lambda^\mu{}_{\nu}x^\nu,
\]
a Dirac spinor transforms as
\[
    \psi(x)\longrightarrow \psi'(x')=S(\Lambda)\psi(x),
\]
where \(S(\Lambda)\) is a spinor representation matrix. It is defined so that
\[
    S(\Lambda)^{-1}\gamma^\mu S(\Lambda)
    =\Lambda^\mu{}_{\nu}\gamma^\nu.
\]
This identity ensures covariance of the Dirac equation.

Infinitesimally, write
\[
    \Lambda^\mu{}_{\nu}=\delta^\mu{}_{\nu}+\omega^\mu{}_{\nu}+O(\omega^2),
\]
where \(\omega_{\mu\nu}=-\omega_{\nu\mu}\). The spinor representation has the
form
\[
    S(\Lambda)=\exp\left(\frac{1}{2}\omega_{\mu\nu}\Sigma^{\mu\nu}\right),
\]
where
\[
    \Sigma^{\mu\nu}=\frac{1}{4}[\gamma^\mu,\gamma^\nu].
\]
These matrices are the Lorentz generators in the spinor representation.

\begin{remarkbox}
Spinors do not transform like vectors. The Lorentz transformation of a spinor is
implemented by matrices \(S(\Lambda)\) acting on spinor components, while the
spacetime point transforms by \(\Lambda^\mu{}_{\nu}\).
\end{remarkbox}

This distinction is central. A spinor index is not a spacetime vector index. It
belongs to a different representation, even though it is tied to spacetime
Lorentz symmetry.

\section{Dirac Field and Dirac Adjoint}

A Dirac field is a spinor-valued field
\[
    \psi:M\to\mathbb{C}^4.
\]
The ordinary Hermitian conjugate \(\psi^\dagger\) does not by itself give a
Lorentz scalar when multiplied by \(\psi\). The correct relativistic adjoint is
the Dirac adjoint
\[
    \bar\psi=\psi^\dagger\gamma^0.
\]
With this definition,
\[
    \bar\psi\psi
\]
is a Lorentz scalar, while
\[
    \bar\psi\gamma^\mu\psi
\]
is a Lorentz vector.

\begin{definitionbox}
The \emph{Dirac adjoint} of a Dirac spinor \(\psi\) is
\[
    \bar\psi=\psi^\dagger\gamma^0.
\]
It is introduced so that spinor bilinears have simple Lorentz transformation
properties.
\end{definitionbox}

Important bilinears include
\[
    \bar\psi\psi,
\]
which is a scalar, and
\[
    j^\mu=\bar\psi\gamma^\mu\psi,
\]
which is a vector current. More bilinears, involving \(\gamma^5\) and
commutators of gamma matrices, appear in advanced treatments of spinor field
theory.

\section{Dirac Lagrangian}

The free Dirac Lagrangian density is
\[
    \mathcal{L}_{\mathrm{D}}
    =\bar\psi(i\gamma^\mu\partial_\mu-m)\psi.
\]
Equivalently,
\[
    \mathcal{L}_{\mathrm{D}}
    =i\bar\psi\gamma^\mu\partial_\mu\psi-m\bar\psi\psi.
\]
This Lagrangian is first order in derivatives. This is one of the main
differences between the Dirac field and the scalar field. The scalar Lagrangian
contains quadratic derivatives; the Dirac Lagrangian contains one derivative.

For a real action one often uses the symmetrized form
\[
    \mathcal{L}_{\mathrm{D}}
    =\frac{i}{2}\left(
        \bar\psi\gamma^\mu\partial_\mu\psi
        -\partial_\mu\bar\psi\gamma^\mu\psi
    \right)-m\bar\psi\psi.
\]
The two forms differ by a total derivative under suitable assumptions. For
classical field equations, the unsymmetrized form is usually sufficient.

\begin{definitionbox}
The free Dirac Lagrangian is
\[
    \mathcal{L}_{\mathrm{D}}=\bar\psi(i\gamma^\mu\partial_\mu-m)\psi.
\]
Its Euler--Lagrange equation is the Dirac equation.
\end{definitionbox}

The mass term
\[
    m\bar\psi\psi
\]
couples left- and right-handed components of the spinor. This becomes especially
important when discussing chirality.

\section{Dirac Equation from the Variational Principle}

In the variational principle, \(\psi\) and \(\bar\psi\) are treated as
independent fields. Varying the action with respect to \(\bar\psi\) gives
\[
    \delta_{\bar\psi}S
    =\int d^{d+1}x\,\delta\bar\psi(i\gamma^\mu\partial_\mu-m)\psi.
\]
Therefore the field equation is
\[
    (i\gamma^\mu\partial_\mu-m)\psi=0.
\]
This is the Dirac equation.

Varying with respect to \(\psi\) gives the adjoint Dirac equation. Starting from
\[
    \mathcal{L}_{\mathrm{D}}
    =i\bar\psi\gamma^\mu\partial_\mu\psi-m\bar\psi\psi,
\]
the variation with respect to \(\psi\) gives, after integration by parts,
\[
    i\partial_\mu\bar\psi\gamma^\mu+m\bar\psi=0.
\]
Equivalently,
\[
    \bar\psi(i\overleftarrow{\partial}_\mu\gamma^\mu+m)=0,
\]
where the arrow indicates that the derivative acts to the left.

\begin{theorembox}
The free Dirac field satisfies
\[
    (i\gamma^\mu\partial_\mu-m)\psi=0,
\]
while the adjoint field satisfies
\[
    i\partial_\mu\bar\psi\gamma^\mu+m\bar\psi=0.
\]
\end{theorembox}

Applying the conjugate first-order operator shows that the Dirac equation implies
the Klein--Gordon equation for each component:
\[
    (\Box+m^2)\psi=0.
\]
Thus the Dirac equation is a relativistic wave equation whose square gives the
mass-shell condition.

\section{Conserved Dirac Current}

The free Dirac Lagrangian is invariant under the global phase transformation
\[
    \psi\longrightarrow e^{i\alpha}\psi,
    \qquad
    \bar\psi\longrightarrow \bar\psi e^{-i\alpha},
\]
where \(\alpha\) is constant. Noether's theorem gives the conserved current
\[
    j^\mu=\bar\psi\gamma^\mu\psi.
\]
Let us verify conservation directly. The Dirac equation is
\[
    i\gamma^\mu\partial_\mu\psi-m\psi=0,
\]
and the adjoint equation is
\[
    i\partial_\mu\bar\psi\gamma^\mu+m\bar\psi=0.
\]
Now compute
\[
    \partial_\mu j^\mu
    =\partial_\mu(\bar\psi\gamma^\mu\psi)
    =(\partial_\mu\bar\psi)\gamma^\mu\psi
     +\bar\psi\gamma^\mu\partial_\mu\psi.
\]
Using the field equations, the two terms cancel, giving
\[
    \partial_\mu j^\mu=0.
\]

\begin{definitionbox}
The Dirac current is
\[
    j^\mu=\bar\psi\gamma^\mu\psi.
\]
It is conserved on shell:
\[
    \partial_\mu j^\mu=0.
\]
\end{definitionbox}

The time component is
\[
    j^0=\bar\psi\gamma^0\psi=\psi^\dagger\psi.
\]
This is nonnegative. This feature was historically important: the Dirac equation
has a positive density associated with its conserved current, unlike the naive
probability interpretation of the Klein--Gordon current.

\section{Coupling to Electromagnetism}

A charged Dirac field couples to electromagnetism by replacing the ordinary
derivative with a gauge covariant derivative. With the convention
\[
    D_\mu=\partial_\mu+iqA_\mu,
\]
the minimally coupled Dirac Lagrangian is
\[
    \mathcal{L}
    =\bar\psi(i\gamma^\mu D_\mu-m)\psi.
\]
Under a local \(U(1)\) transformation,
\[
    \psi\longrightarrow e^{i\alpha(x)}\psi,
\]
the gauge potential transforms as
\[
    A_\mu\longrightarrow A_\mu-\frac{1}{q}\partial_\mu\alpha.
\]
Then
\[
    D_\mu\psi\longrightarrow e^{i\alpha(x)}D_\mu\psi.
\]
Thus the Lagrangian is locally gauge invariant.

If the electromagnetic field is dynamical, the full Lagrangian is
\[
    \mathcal{L}
    =-\frac{1}{4}F_{\mu\nu}F^{\mu\nu}
     +\bar\psi(i\gamma^\mu D_\mu-m)\psi.
\]
The Dirac equation becomes
\[
    (i\gamma^\mu D_\mu-m)\psi=0.
\]
Varying the action with respect to \(A_\mu\) gives Maxwell's equation with Dirac
matter current. With the above convention, the electric current is proportional
to
\[
    j^\mu=q\bar\psi\gamma^\mu\psi,
\]
up to the sign convention used in the interaction term.

\begin{remarkbox}
Minimal electromagnetic coupling has the same logic for scalar and spinor
matter: local phase symmetry forces the introduction of the gauge potential and
the covariant derivative.
\end{remarkbox}

The Dirac field is therefore the simplest classical relativistic model of
charged spin-one-half matter coupled to electromagnetism.

\section{Chiral Decomposition}

In four spacetime dimensions one defines
\[
    \gamma^5=i\gamma^0\gamma^1\gamma^2\gamma^3.
\]
It satisfies
\[
    (\gamma^5)^2=\mathbb{I},
    \qquad
    \{\gamma^5,\gamma^\mu\}=0.
\]
The chiral projection operators are
\[
    P_L=\frac{1}{2}(\mathbb{I}-\gamma^5),
    \qquad
    P_R=\frac{1}{2}(\mathbb{I}+\gamma^5).
\]
They obey
\[
    P_L^2=P_L,
    \qquad
    P_R^2=P_R,
    \qquad
    P_LP_R=0,
    \qquad
    P_L+P_R=\mathbb{I}.
\]
The left- and right-handed parts of a Dirac spinor are
\[
    \psi_L=P_L\psi,
    \qquad
    \psi_R=P_R\psi.
\]

For a massless Dirac field, the Dirac equation separates into independent left-
and right-handed equations. The mass term couples them:
\[
    m\bar\psi\psi=m(\bar\psi_L\psi_R+\bar\psi_R\psi_L).
\]
Thus chirality is especially important for massless or approximately massless
fermions.

\begin{definitionbox}
Chiral projectors split a Dirac spinor into left- and right-handed components:
\[
    \psi=\psi_L+\psi_R,
    \qquad
    \psi_L=P_L\psi,
    \qquad
    \psi_R=P_R\psi.
\]
\end{definitionbox}

Chirality is not just a mathematical refinement. It is essential in the weak
interaction, where left- and right-handed fermions transform differently under
the gauge group.

\section{Weyl and Majorana Spinors}

A Weyl spinor is a spinor of definite chirality. In four dimensions, it is
obtained by imposing
\[
    \psi=\psi_L
\]
or
\[
    \psi=\psi_R.
\]
For massless fields, left- and right-handed Weyl spinors can be treated
independently. A mass term generally couples opposite chiralities and therefore
cannot be written for a single isolated Weyl spinor in the same way as for a
Dirac spinor.

A Majorana spinor satisfies a reality condition relating the spinor to its charge
conjugate. Schematically,
\[
    \psi^c=\psi.
\]
The precise form of charge conjugation depends on the gamma-matrix convention.
Majorana spinors reduce the number of independent real degrees of freedom.

\begin{remarkbox}
Dirac, Weyl, and Majorana spinors are not three unrelated objects. They are
different ways of organizing spin-one-half degrees of freedom, depending on
chirality, reality conditions, dimension, signature, and gauge representation.
\end{remarkbox}

In many physical applications, whether a Majorana condition is possible depends
on the gauge charges of the field. A charged field generally cannot be Majorana
with respect to the same \(U(1)\) charge, because charge conjugation reverses the
charge.

\section{Classical Spinor Fields vs Quantum Fermions}

The Dirac equation can be studied as a classical field equation. One may solve
it, analyze its currents, couple it to gauge fields, and study its Lorentz
covariance. However, the physical interpretation of spinor fields is deeply
quantum.

In quantum field theory, spin-one-half fields are fermionic. Their field
operators satisfy anticommutation relations, not commutation relations. This is
connected with the spin-statistics theorem. A fully consistent quantum treatment
of the Dirac field therefore requires fermionic quantization.

In classical variational field theory, one often treats spinor components as
Grassmann-valued variables. Grassmann variables anticommute:
\[
    \theta_1\theta_2=-\theta_2\theta_1.
\]
This formalism anticipates the algebraic structure needed for fermionic quantum
fields and path integrals.

\begin{definitionbox}
A \emph{Grassmann variable} is an anticommuting variable. Classical fermionic
field theory is often formulated using Grassmann-valued spinor fields, especially
when preparing for quantization.
\end{definitionbox}

For a first course in classical field theory, it is acceptable to first treat the
Dirac equation formally as a classical relativistic wave equation. But one should
not confuse this formal classical description with the full physics of fermions.

\section{What Is Classical and What Is Already Anticipating QFT?}

Several parts of the Dirac theory are genuinely classical field-theoretic:
\begin{enumerate}
    \item the Lorentz covariance of the Dirac equation;
    \item the construction of gamma matrices and spinor representations;
    \item the variational principle for the Dirac action;
    \item the conserved current associated with global phase symmetry;
    \item minimal coupling to gauge fields.
\end{enumerate}

Other parts point beyond ordinary classical field theory:
\begin{enumerate}
    \item the interpretation of \(\psi^\dagger\psi\) as probability density for a
    single relativistic particle;
    \item the existence of antiparticles;
    \item fermionic anticommutation relations;
    \item the spin-statistics connection;
    \item Grassmann path integrals.
\end{enumerate}

The safest interpretation in this course is the following. The Dirac field is a
classical field variable with a Lorentz-covariant first-order action. Its full
physical meaning as a fermion field emerges only in quantum field theory.

\begin{summarybox}
Spinor fields carry spin representations of the Lorentz group. The gamma
matrices satisfy the Clifford algebra
\(\{\gamma^\mu,\gamma^\nu\}=2\eta^{\mu\nu}\mathbb{I}\), and the Dirac equation is
\((i\gamma^\mu\partial_\mu-m)\psi=0\). The Dirac adjoint
\(\bar\psi=\psi^\dagger\gamma^0\) allows Lorentz-covariant bilinears, including
the conserved current \(j^\mu=\bar\psi\gamma^\mu\psi\). Minimal coupling replaces
\(\partial_\mu\) by \(D_\mu\), and chiral projectors split a Dirac spinor into
left- and right-handed components. The classical formalism is meaningful, but
the physical interpretation of fermions belongs to quantum field theory.
\end{summarybox}
